% !TEX root =  mosaic.tex
\section{Metadata Extraction and GPU Profiling}
\label{sec:metadata}

This section explains how \name converts a transformed program and a target \gpu into the constants required by \autoref{eq:two-level}. \name instantiates its analytical model from two sources: compile-time metadata extracted from the transformed program and target-specific constants obtained once per \gpu architecture. The first source describes the computation and its transformed loop structure; the second describes the resource limits and effective costs of the target. Together, they determine the candidate variables, resource footprints, feasibility constraints, and cost terms used by the solver.

\subsection{Compile-Time Metadata}
\label{sec:compile-time-metadata}

The transformed representation exposes the loop order and mapping decisions the scheduler already made, so \name optimizes execution space parameters without redoing them.

The metadata falls into several groups tied to the model. These groups are chosen to capture the information the cost model requires: loop structure drives candidate generation, array metadata drives traffic and placement decisions, statement metadata drives compute and synchronization cost, and GPU-specific mapping metadata closes the occupancy and register loop. Loop structure records dimension names, bounds, schedule order, and step sizes; with \(N_d\) the extent of dimension \(d\), these define the tile and block candidate sets. Dimension annotations record permutable and coincident (i.e., parallelizable) dimensions, from which \name classifies each loop as parallel (\(\mathcal{P}\), mappable to independent blocks/threads), reduction (\(\mathcal{R}\), absent from every write footprint), or serial (\(\mathcal{D}\), neither of the former); the formal definition is \autoref{eq:dimension-roles} in \autoref{sec:abstraction}. Array metadata records affine accesses, read/write polarity, layout, coalescing direction, and tile footprints, determining memory traffic, shared memory capacity, accumulator placement, and staging candidates. Statement metadata records arithmetic counts, memory-reference counts, and synchronization sites, which instantiate compute and barrier costs. GPU-specific mapping metadata adds block/thread mappings, shared memory staging candidates, register-blocking factors \(\rblk_d=\tile_d/b_d\), and coalescing.

The shared memory placement candidates are restricted to read-only references whose tile footprints are functions of the candidate tile sizes; \autoref{tab:placement-rules} assigns every other reference a fixed memory space. This restriction keeps the outer enumeration small while retaining the placements that materially affect reuse, shared memory footprint, and occupancy.

\begin{table}[htbp]
  \caption{Placement rules.}
  \label{tab:placement-rules}
  \centering
  \scriptsize
  \setlength{\tabcolsep}{0pt}
  \renewcommand{\arraystretch}{1.05}

  \begin{tabularx}{\columnwidth}{
    @{}
    >{\raggedright\arraybackslash}X
    @{}
  }
    \toprule
    \textbf{Array polarity and placement decision} \\
    \midrule

    \textbf{Read/write accumulators.}\par
    Place in registers to keep private accumulated outputs
    close to the thread. \\

    \addlinespace[2.5pt]

    \textbf{Write-only coalesced references.}\par
    Keep in global memory and stream the output directly. \\

    \addlinespace[2.5pt]

    \textbf{Write-only uncoalesced or transposed references.}\par
    Use shared memory to restore coalescing before global
    writeback. \\

    \addlinespace[2.5pt]

    \textbf{Read-only references.}\par
    Form \(\Sigma_{\mathit{cand}}\); the cost model decides
    whether staging is profitable. \\

    \bottomrule
  \end{tabularx}
\end{table}

\autoref{fig:gemm-schedule-metadata} illustrates metadata construction using the default PPCG schedule of the running GEMM example and the dimension sets defined in \autoref{eq:dimension-roles}. In general, \name requires only the categories above; polyhedral compilers expose them through schedule trees and affine access maps, while TVM-style compilers expose them through schedule traces and tensor access metadata.

% Figure palette
\definecolor{color1}{RGB}{199,56,73}
\definecolor{color2}{RGB}{248,189,211}
\definecolor{color3}{RGB}{191,170,210}
\definecolor{color4}{RGB}{147,197,235}
\definecolor{color5}{RGB}{157,159,163}

\definecolor{color1}{RGB}{199,56,73}
\definecolor{color2}{RGB}{248,189,211}
\definecolor{color3}{RGB}{191,170,210}
\definecolor{color4}{RGB}{147,197,235}
\definecolor{color5}{RGB}{157,159,163}

\begin{figure}[htbp]
\centering
% \setlength{\abovecaptionskip}{10pt}
\resizebox{0.96\columnwidth}{!}{%
\begin{tikzpicture}[
  font=\scriptsize,
  box/.style={
    draw=color5!82!black,
    fill=white,
    rounded corners=2pt,
    inner sep=2.6pt,
    text width=#1,
    align=left
  },
  textonly/.style={
    text width=#1,
    align=left,
    inner sep=0pt
  },
  hdr/.style={
    draw=color5!82!black,
    rounded corners=2pt,
    inner sep=2.5pt,
    fill=color5!12,
    font=\scriptsize\bfseries,
    align=center
  },
  arr/.style={
    -{Latex[length=1.55mm]},
    line width=0.55pt,
    draw=color1!85!black
  },
  grpborder/.style={
    draw=color5!72,
    rounded corners=3pt,
    line width=0.45pt,
    dashed
  },
]

\node[
  hdr,
  text width=72mm,
  fill=color4!30,
  draw=color4!62!black
] (src_hdr) at (0,0)
  {Schedule metadata (from any compiler)};

\node[
  textonly=41mm,
  anchor=north west
] (bandone_text)
  at ([xshift=-35.8mm,yshift=-2.2mm]src_hdr.south)
{%
  \(\triangleright\) \textbf{band\(_1\)}\\
  \quad permutable: 1\\
  \quad parallelizable: 1\\
  \quad dims: \(i,j\)\\
  \quad bounds: \(0\le i<N_i\), \(0\le j<N_j\)\\
  \quad tiles: \(t_i,t_j\)\\
  \quad \(1\le t_i<N_i\), \(1\le t_j<N_j\)\\
  \quad blocks: \(b_i,b_j\)\\
  \quad \(1\le b_i\le t_i\), \(1\le b_j\le t_j\)
};

\node[
  box=25mm,
  dashed,
  anchor=north east
] (bandtwo_box)
  at ([xshift=35.8mm,yshift=-2.2mm]src_hdr.south)
{%
  \(\triangleright\) \textbf{band\(_2\)}\\[-0.2pt]
  \quad permutable: 1\\
  \quad parallelizable: 0\\
  \quad map: sequential\\
  \quad dim: \(k\)\\
  \quad bound: \(0\le k<N_k\)\\
  \quad tile: \(t_k\); \(1\le t_k<N_k\)
};

\node[
  textonly=68mm,
  anchor=north west
] (mapping_text)
  at ([xshift=0mm,yshift=-1.5mm]bandone_text.south west)
{%
  \quad block map: \(i\to b_0,\ j\to b_1\);\quad
  thread map: \(i\to t_0,\ j\to t_1\)
};

\node[
  textonly=68mm,
  anchor=north west
] (access_text)
  at ([xshift=0mm,yshift=-1.15mm]mapping_text.south west)
{%
  \makebox[\linewidth][l]{%
    \textbf{Array access:}\quad
    Read-only: \(A[i,k],B[k,j]\)\hfill
    Read-write: \(C[i,j]\)}
};

\begin{scope}[on background layer]
  \node[
    draw=color4!55!black,
    fill=color4!5,
    rounded corners=2pt,
    inner sep=2.6pt,
    fit=(bandone_text)(bandtwo_box)(mapping_text)(access_text)
  ] (loop_box) {};
\end{scope}

\node[
  box=72mm,
  draw=color1!68!black,
  fill=color2!13,
  anchor=north,
  inner xsep=2.8pt,
  inner ysep=2.5pt
] (code_box)
  at ([yshift=-1.1mm]loop_box.south)
{%
  \textbf{Scheduled GEMM instance:}\\[-0.2pt]
  {\ttfamily\fontsize{4.75pt}{5.75pt}\selectfont
  \begin{tabular}{@{}l@{\hspace{1.15em}}l@{}}
  for (it, jt) in parallel tiles:
    & // band$_1$ $\to$ blocks/threads\\
  \quad for (i, j) in tile(it, jt):
    & \\
  \qquad C[i][j] *= beta;
    & // $S_0$: 1 mul\\
  \quad for kt in 0..N\_k step t\_k:
    & // band$_2$ $\to$ sequential\\
  \qquad for (i, j) in tile, k in kt-tile:
    & \\
  \qquad\quad C[i][j] += alpha*A[i][k]*B[k][j];
    & // $S_1$: 2 mul + 1 add
  \end{tabular}}
};

\coordinate (lower_top)
  at ([yshift=-3.0mm]code_box.south);
\node[
  hdr,
  text width=28mm,
  fill=color4!24,
  draw=color4!62!black,
  anchor=north west
] (hw_hdr)
  at ([xshift=0mm]src_hdr.west |- lower_top)
  {Hardware profile};

\node[
  box=28mm,
  draw=color4!58!black,
  fill=color4!6,
  anchor=north,
  yshift=-1.0mm,
  minimum height=12.0mm,
  inner ysep=1.5pt
] (hw_box)
  at (hw_hdr.south)
{%
  {\fontsize{5.15pt}{6.15pt}\selectfont
  \textbf{Capacities}\\
  \(R_{\mathrm{SM}},
    S_{\mathrm{SM}},
    S_{\mathrm{block}},
    B_{\mathrm{SM}},
    T_{\mathrm{SM}}\)\\
  \(N_{\mathrm{SM}},
    T_{\max},
    w_{\mathrm{warp}},
    g_{\mathrm{coal}},
    \ldots\)\\
  \textbf{Effective costs}\\
  \(\Lambda_g,
    \Lambda_s,
    \Lambda_r,
    \Lambda_f,
    \Lambda_y\)}
};

\begin{scope}[on background layer]
  \node[
    grpborder,
    inner sep=1.8pt,
    fit=(hw_hdr)(hw_box)
  ] (hw_grp) {};
\end{scope}

\node[
  hdr,
  text width=37.5mm,
  fill=color3!32,
  draw=color3!68!black,
  anchor=north east
] (out_hdr)
  at ([xshift=0mm]src_hdr.east |- lower_top)
  {\name analytical metadata};

\node[
  box=37.5mm,
  draw=color3!64!black,
  fill=color2!9,
  anchor=north,
  yshift=-1.0mm,
  minimum height=12.0mm,
  inner ysep=1.5pt
] (out_box)
  at (out_hdr.south)
{%
  {\fontsize{5.15pt}{6.15pt}\selectfont
  \(\mathcal P=\{i,j\},\
    \mathcal R=\{k\},\
    \mathcal D=\emptyset\)\\
  Tile vars: \(t_i,t_j,t_k\);\quad
  Block vars: \(b_i,b_j\)\\
  Reg.: \(r_d=t_d/b_d\ (d\in\mathcal P)\)\\
  {\fontsize{4.55pt}{6.15pt}\selectfont
  \(\Sigma_{\mathrm{cand}}=\{A,B\};\quad
    \Sigma=\{\emptyset,\{A\},\{B\},\{A,B\}\}\)}\\
  Accumulator in registers: \(C\)}
};

\begin{scope}[on background layer]
  \node[
    grpborder,
    inner sep=1.8pt,
    fit=(out_hdr)(out_box)
  ] (out_grp) {};
\end{scope}

\begin{scope}[on background layer]
  \node[
    grpborder,
    inner sep=1.8pt,
    fit=(src_hdr)(loop_box)(code_box)
  ] (sched_grp) {};
\end{scope}

\coordinate (sched_arrow_start)
  at (out_grp.north |- sched_grp.south);

\draw[arr]
  (sched_arrow_start)
  -- (out_grp.north);

\coordinate (hw_exit)
  at (hw_grp.east |- hw_box.center);

\coordinate (hw_entry)
  at (out_grp.west |- hw_box.center);

\draw[arr]
  (hw_exit)
  -- (hw_entry);

\end{tikzpicture}%
}

\caption{
  GEMM metadata construction.
  Schedule metadata from the PPCG schedule tree, combined with a
  one-time hardware profile, yields the \name analytical metadata,
  with the dimension sets following \autoref{eq:dimension-roles}.
}

\Description{
  A compact single-column flow diagram in which a rectangular
  schedule-metadata panel contains a GEMM schedule-tree instance,
  while a one-time hardware profile and the schedule metadata feed
  \name analytical metadata.
}

\label{fig:gemm-schedule-metadata}
\end{figure}

In this example, the dimension sets defined in \autoref{eq:dimension-roles} give \(\mathcal{P}=\{i,j\}\), which creates \(\tile_i,\tile_j\) and \(b_i,b_j\); \(\mathcal{R}=\{k\}\), which creates \(\tile_k\) but no block dimension. The register-blocking factors \(\rblk_i=\tile_i/b_i\) and \(\rblk_j=\tile_j/b_j\) follow from these tile and block choices rather than entering as separate variables. The read-only operands give \(\Sigma_{\mathit{cand}}=\{A,B\}\), whose four binary-based placement schemes change the footprint, traffic, and resident-block bound used in \autoref{sec:model}.

\subsection{GPU Target Profiling}
\label{sec:gpu-target-profiling}

The target description combines hardware constants queried from the device with effective costs measured once per architecture and reused across all kernels for that target. The hardware constants, including SM count, warp size, resident block and thread limits, shared memory capacities, register-file size, and block-size limit (\autoref{tab:model-notation}), bound the resident-block count, occupancy, and feasible domain. The effective costs, the \(\Lambda\)-family constants of \autoref{tab:model-notation}, come from targeted microbenchmarks and enter the matching terms of \(T_{\mathit{block}}\). They serve relative ranking and analytical selection, not cycle-accurate prediction; the model captures the resource interactions that shape the execution space landscape and leaves backend scheduling, caching, and pipeline overlap as lower-level effects.

\begin{table}[htbp]
  \caption{Model notation.}
  \label{tab:model-notation}
  \centering
  \scriptsize
  \setlength{\tabcolsep}{3pt}
  \begin{tabular}{p{0.30\linewidth}p{0.59\linewidth}}
    \toprule
    Symbol & Meaning \\
    \midrule
    \multicolumn{2}{l}{\textbf{Decision variables}} \\[1pt]
    \(\tile_d,\ b_d,\ \rblk_d\) & Tile size, thread-block size, and per-thread
                              register-blocking factor for dimension~\(d\). \\
    \(\work,\ S \in \Sigma\)   & Per-thread output count; shared memory
                              placement scheme (subset of \(\Sigma_{\mathit{cand}}\)). \\
    \midrule
    \multicolumn{2}{l}{\textbf{Architecture constants}
      \textit{(queried or profiled once per GPU)}} \\[1pt]
    \(\Nsm,\Bsm,\Tsm,\Tmax\) &
      SM count; max resident blocks and threads per SM; max threads per block. \\
    \(\Ssm,\Sblock,\Rsm\) &
      Per-SM and per-block shared memory capacity; SM register-file size. \\
    \(\Lg,\Ls,\Lf,\Ly\) &
      Effective costs for global access, shared-memory access, FMA,
      and barrier synchronization. \\
    \(\Lr,\ \gcoal\) &
      Register-access cost; coalescing granularity for loads. \\
    \midrule
    \multicolumn{2}{l}{\textbf{Key derived quantities}} \\[1pt]
    \(\phi_{\square}\) &
      Reads-per-FMA ratio for placement class
      \(\square\!\in\!\{\mathit{smem},\mathit{gmem}\}\);
      links register blocking to memory traffic in \(T_{\mathit{compute}}\). \\
    \(\operatorname{smem}(\tvec;S),\ |A|_{\tvec}\) &
      Shared-memory footprint and operand tile footprint. \\
    \(K_{\mathit{outer}}\) &
      Cooperative load rounds \(\prod_{r\in\mathcal{R}}\lceil N_r/\tile_r\rceil\)
      over tiled reduction dimensions. \\
    \(bps,\ R_{\mathit{avail}},\ \hat{R}\) &
      Resident blocks per SM; available registers per thread;
      conservative register working-set estimate. \\
    \bottomrule
  \end{tabular}
\end{table}

In practice, the microbenchmark suite is fewer than ten short kernels, each measuring one cost constant in isolation, and completes in under two minutes on any of the three GPUs. The profiled constants depend only on the architecture, not on any kernel or problem size, so this one-time cost amortizes across all kernels for that target.
